\documentclass[journal,10pt,twocolumn]{IEEEtran}
\usepackage{amsmath,amssymb,amsfonts}
\usepackage{cite}
\usepackage{graphicx}
\usepackage{booktabs}
\usepackage{fontspec}
\usepackage[english, bidi=basic, provide=*]{babel}
\babelprovide[import, onchar=ids fonts]{english}
\babelfont{rm}{Noto Sans}
\usepackage{hyperref}

\begin{document}

\title{Operation Dim Star: Photometric Exoplanet Detection, Orbital Characterization, and Habitability Assessment}

\author{Alliance~Astronomical~Analysis~Squad,~\IEEEmembership{Exoplanet~Intelligence~Unit}%
\thanks{Submission for STAC WARS Technical Evaluation Task.}%
\thanks{Special Technical Operations Command (STAC) - Document Ref: IEEE-STAC-2026-03.}}

\markboth{IEEE Transactions on Exoplanetary Intelligence,~Vol.~14,~No.~8,~October~2026}%
{Alliance Squad: Operation Dim Star Technical Analysis Report}

\maketitle

\begin{abstract}
Photometric light curve data for three target stellar systems (STAC-1, STAC-2, and STAC-3) were analyzed to locate viable terrestrial worlds for strategic base relocation. Utilizing transit photometry and Kepler's Third Law, planetary radii, orbital periods, semi-major axes, and circumstellar habitable zone parameters were calculated. The analysis detected three distinct exoplanet candidates: STAC-1b ($R_p = 5.98\,R_\oplus$, $P = 4.0\text{d}$), STAC-2b ($R_p = 2.76\,R_\oplus$, $P = 80.0\text{d}$), and STAC-3b ($R_p = 1.64\,R_\oplus$, $P = 40.0\text{d}$). While STAC-1b and STAC-2b are non-viable due to high thermal irradiance and gas-giant scale, STAC-3b resides closest to its host star's habitable zone ($d_{\text{HZ}} \approx 0.283\text{ AU}$). STAC-3b is concluded to be the most promising target for immediate astronomical follow-up and site settlement.
\end{abstract}

\begin{IEEEkeywords}
Exoplanets, Transit Photometry, Kepler's Third Law, Planetary Radii, Habitable Zone, Super-Earth.
\end{IEEEkeywords}

\section{Introduction}
\IEEEPARstart{F}{inding} habitable terrestrial worlds beyond charted Imperial reach requires systematic analysis of stellar photometric variations. Spaceborne telescopic monitoring of targets STAC-1, STAC-2, and STAC-3 revealed periodic flux reductions indicative of transiting exoplanets. This paper outlines the quantitative mathematical derivations used to compute planetary physical dimensions, orbital mechanics, and habitable zone boundaries for all three candidate systems.

\section{Detection and Light Curve Analysis}
When an exoplanet transits its host star along our line of sight, it blocks a portion of the stellar disk proportional to its cross-sectional area, reducing observed stellar flux. The fractional transit depth $\delta$ is given by:
\begin{equation}
\delta = \frac{F_{\text{normal}} - F_{\text{transit}}}{F_{\text{normal}}}
\end{equation}

\subsection{STAC-1 System}
Photometric observations indicate flux dips from $1.0000$ to $0.9970$ at times $t = 3, 8, 12, 16, 20\text{ days}$. Following initial transit at $t = 3\text{d}$, periodic recurrence occurs consistently every $\Delta t = 4.0\text{ days}$.
\begin{equation}
P_1 = 4.0\text{ days} = 0.01095\text{ years}
\end{equation}
\begin{equation}
\delta_1 = \frac{1.0000 - 0.9970}{1.0000} = 0.0030 \quad (0.30\%)
\end{equation}

\subsection{STAC-2 System}
Light curve observations spanning 200 days show flux drops to $0.9990$ at $t = 20, 100, 180\text{ days}$. The recurrence interval is strictly uniform at $\Delta t = 80.0\text{ days}$.
\begin{equation}
P_2 = 80.0\text{ days} = 0.21903\text{ years}
\end{equation}
\begin{equation}
\delta_2 = \frac{1.0000 - 0.9990}{1.0000} = 0.0010 \quad (0.10\%)
\end{equation}

\subsection{STAC-3 System}
STAC-3 exhibits recurring flux dips at $t = 20, 60, 100, 140, 170\text{ days}$ with mean flux during transit $F_{\text{transit, avg}} \approx 0.9991$. Transit periodicity recurs every $\Delta t = 40.0\text{ days}$.
\begin{equation}
P_3 = 40.0\text{ days} = 0.10951\text{ years}
\end{equation}
\begin{equation}
\delta_3 = \frac{1.0000 - 0.9991}{1.0000} = 0.0009 \quad (0.09\%)
\end{equation}

\section{Planetary Radius Calculations}
Using the geometric relation between transit depth and planetary radius:
\begin{equation}
\delta \approx \left(\frac{R_p}{R_s}\right)^2 \implies R_p = R_s \cdot \sqrt{\delta}
\end{equation}
Given standard constants: $R_\odot = 696,000\text{ km}$ and $R_\oplus = 6,371\text{ km}$.

\subsection{STAC-1b Radius}
$R_{s1} = 1.00 R_\odot = 696,000\text{ km}$.
\begin{equation}
R_{p1} = 696,000 \cdot \sqrt{0.0030} = 38,121.5\text{ km}
\end{equation}
\begin{equation}
R_{p1} = \frac{38,121.5}{6,371} = 5.98\,R_\oplus \quad \text{(Neptune-sized)}
\end{equation}

\subsection{STAC-2b Radius}
$R_{s2} = 0.80 R_\odot = 556,800\text{ km}$.
\begin{equation}
R_{p2} = 556,800 \cdot \sqrt{0.0010} = 17,607.7\text{ km}
\end{equation}
\begin{equation}
R_{p2} = \frac{17,607.7}{6,371} = 2.76\,R_\oplus \quad \text{(Super-Earth)}
\end{equation}

\subsection{STAC-3b Radius}
$R_{s3} = 0.50 R_\odot = 348,000\text{ km}$.
\begin{equation}
R_{p3} = 348,000 \cdot \sqrt{0.0009} = 10,440.0\text{ km}
\end{equation}
\begin{equation}
R_{p3} = \frac{10,440.0}{6,371} = 1.64\,R_\oplus \quad \text{(Terrestrial Super-Earth)}
\end{equation}

\section{Orbital Distance Computations}
Applying Kepler's Third Law in canonical units ($P$ in years, $M$ in $M_\odot$, $a$ in AU):
\begin{equation}
P^2 = \frac{a^3}{M} \implies a = \left(P^2 \cdot M\right)^{1/3}
\end{equation}
where $1\text{ AU} = 149,600,000\text{ km}$.

\subsection{STAC-1b Orbital Distance}
\begin{equation}
a_1 = \left[(0.01095)^2 \cdot 1.00\right]^{1/3} = 0.0493\text{ AU}
\end{equation}
$a_1 = 0.0493 \times 149,600,000 = 7,375,280\text{ km}$.

\subsection{STAC-2b Orbital Distance}
\begin{equation}
a_2 = \left[(0.21903)^2 \cdot 0.80\right]^{1/3} = 0.3373\text{ AU}
\end{equation}
$a_2 = 0.3373 \times 149,600,000 = 50,460,080\text{ km}$.

\subsection{STAC-3b Orbital Distance}
\begin{equation}
a_3 = \left[(0.10951)^2 \cdot 0.50\right]^{1/3} = 0.1817\text{ AU}
\end{equation}
$a_3 = 0.1817 \times 149,600,000 = 27,182,320\text{ km}$.

\section{Habitability Zone Evaluation \& Selection}
The theoretical circumstellar habitable zone boundary $d_{\text{HZ}}$ is estimated as:
\begin{equation}
d_{\text{HZ}} \approx \sqrt{\frac{L}{L_\odot}}\text{ AU}
\end{equation}

\begin{itemize}
    \item \textbf{STAC-1 System ($L_1 = 1.00\,L_\odot$):} $d_{\text{HZ1}} = 1.000\text{ AU}$. STAC-1b ($a_1 = 0.0493\text{ AU}$) is heavily overheated.
    \item \textbf{STAC-2 System ($L_2 = 0.40\,L_\odot$):} $d_{\text{HZ2}} = 0.632\text{ AU}$. STAC-2b ($a_2 = 0.3373\text{ AU}$) orbits interior to the habitable zone edge.
    \item \textbf{STAC-3 System ($L_3 = 0.08\,L_\odot$):} $d_{\text{HZ3}} = 0.283\text{ AU}$. STAC-3b ($a_3 = 0.1817\text{ AU}$) orbits closest to its stellar habitable region.
\end{itemize}

\begin{table*}[t]
\caption{System Parameter Summary and Exoplanetary Comparisons}
\label{tab:summary}
\centering
\begin{tabular}{l c c c}
\toprule
\textbf{Parameter} & \textbf{STAC-1b} & \textbf{STAC-2b} & \textbf{STAC-3b (Selected Target)} \\
\midrule
Orbital Period ($P$) & $4.00\text{ days}$ & $80.00\text{ days}$ & $\mathbf{40.00\text{ days}}$ \\
Transit Depth ($\delta$) & $0.0030\ (0.30\%)$ & $0.0010\ (0.10\%)$ & $\mathbf{0.0009\ (0.09\%)}$ \\
Planetary Radius ($R_p$) & $38,121.5\text{ km}\ (5.98\,R_\oplus)$ & $17,607.7\text{ km}\ (2.76\,R_\oplus)$ & $\mathbf{10,440.0\text{ km}\ (1.64\,R_\oplus)}$ \\
Classification & Neptune-sized Gas Planet & Large Super-Earth & $\mathbf{Terrestrial\ Super\text{-}Earth}$ \\
Orbital Distance ($a$) & $0.0493\text{ AU}\ (7.38 \times 10^6\text{ km})$ & $0.3373\text{ AU}\ (50.46 \times 10^6\text{ km})$ & $\mathbf{0.1817\text{ AU}\ (27.18 \times 10^6\text{ km})}$ \\
Habitable Zone ($d_{\text{HZ}}$) & $1.000\text{ AU}$ & $0.632\text{ AU}$ & $\mathbf{0.283\text{ AU}}$ \\
Final Decision & Reject (Gas Giant) & Reject (Too Hot) & $\mathbf{ACCEPTED\ FOR\ BASE\ LOCATION}$ \\
\bottomrule
\end{tabular}
\end{table*}

\section{Conclusion and Recommendation}
\textbf{STAC-3b} is selected as the primary target for further study based on:
\begin{enumerate}
    \item \textbf{Terrestrial Scale ($1.64\,R_\oplus$):} Confirms a rocky surface without a volatile gas envelope.
    \item \textbf{Habitable Proximity:} Best suited thermal regime around its low-luminosity M-dwarf host.
    \item \textbf{Atmospheric Follow-Up:} Transmission spectroscopy during transits is recommended to confirm water vapor ($H_2O$) and oxygen/ozone ($O_3$) absorption lines.
\end{enumerate}

\end{document}